\cleardoublepage{}
\begin{center}
    \bfseries \zihao{3} 致~谢
\end{center}

\par 四年的本科求学时光行将结束，毕业设计的撰写也终于走到尾声。回望这段从课堂到毕设、从基础到方向的过程，我得到了许多老师与同学的帮助，谨在此一并致谢。

\par 首先要诚挚感谢我的指导教师朱霖潮老师。
本课题在选题、文献调研、技术路线、实验设计、论文撰写等各个阶段，朱老师都给予了具体而切中要害的指导。
特别是在审核意见中提出的“强化核心算法创新、补充关键技术模块算法细节”，让我从最初较为工程化的方案重新聚焦为 RATS 与 DAGPlanner 两个可独立消融的算法模块，这一重构是本文工作能够形成体系的关键。
朱老师严谨的学术态度与宽容的指导方式，对我形成独立思考能力与工程克制感有着深远影响，在此向朱老师致以最诚挚的敬意与感谢。

\par 同样要感谢计算机科学与技术学院四年来所有任课老师。
正是各位老师在程序设计、数据结构、操作系统、计算理论、机器学习等课程上的扎实讲授，使我具备了完成这份毕业设计所必需的工程素养与理论基础。

\par 感谢答辩与评审环节的各位老师与盲审专家。
他们对本课题在中期阶段提出的具体修改意见，是本文从“能跑通”走向“可论证”的重要推动力。

\par 感谢实验室与同班同学在本课题进行过程中给予的讨论与互助；
也感谢我所在公司的实习同事在我利用工作之余推进毕设时给予的理解与包容。

\par 最后，要感谢我的家人。
本科四年远在异乡求学，正是家人无条件的支持与信任，让我得以在求学与实习的双重压力下专注完成本课题，并最终在论文中诚实地呈现工作的边界与不足。

\par 智能体与大语言模型领域日新月异，本文工作只是其中很小的一个切片。
这份毕业设计既是本科学习的一个总结，也是我未来继续在该领域深入探索的一个新起点。

\vspace{1em}
\begin{flushright}
彭超琦\\
2026 年 6 月
\end{flushright}
